\documentclass[a4paper]{article}
\usepackage{luatexja}
\usepackage{amsmath}
\usepackage{amssymb}
\usepackage{geometry}
\usepackage{hyperref}

\geometry{margin=25mm}

\title{双対ローター位相遅れとしてのド・ブロイ波：\\
粒子内在双対時空の幾何としての量子力学}

\author{https://github.com/hypernumbernet}
\date{\today}

\begin{document}
\maketitle

\begin{abstract}
二重時空理論におけるすべての質量粒子は双対ローターを運び、それはすでに \(\mathbb{C}^2\) 上の \(\mathrm{SU}(2)\) の元である。そのローターの不変量は三つに分けておかねばならない。指標 \(\operatorname{Tr} R(\beta)=2\cos(\lVert\beta\rVert/2)\) が実であるのは、二つの計算基底スピノル上の反対の半角位相の和だからである。複素行列要素 \(\langle\uparrow\rvert R(\beta)\lvert\uparrow\rangle=\cos(\theta/2)-i n_z\sin(\theta/2)\) は観測者軸では純位相、赤道面では実余弦であり、その絶対値は双対軸の極角傾斜である。ド・ブロイ位相と同定されるスカラー \(\mathrm{U}(1)\) 因子は、それらとは別の第三の量である。観測者と粒子は双対ローター像のヒルベルト対で比べられ、それはそれ自身 \(\mathrm{SU}(2)\) に留まる相対双対ローターの行列要素に等しい。その相対ローターは群の積であり、指標は二つの双対軸の整列を記録する。同じ行列要素は、内部双対回転だけ異なる静止系ディラック解の重なりである。書かれた反対符号の粒子作用は不整合を自由にし、遅れは線形に蓄積しうる。調和模型は同符号の運動項に限り現れ、振動数は \(\sqrt{2m}\) であって \(E/\hbar\) ではない。運動量空間のディラック方程式は \(\mathrm{Cl}(3,1)\) におけるミンコフスキー二次形式の一次因子である。シュレーディンガーおよびクライン–ゴルドンの動力学、収縮、重力と量子非局所性の統一は幾何的な読みとして残る。本稿の位相 \(J\) は振動子の診断量であり、主論文のパラメータ空間キリング二次形式ではない。
\end{abstract}

\section{はじめに}

すべての質量粒子が波長 \(\lambda=h/p\) の波を伴うというド・ブロイの1924年の仮説は、物理学における最も深い経験的事実の一つであり続けている~\cite{debroglie1924}。標準量子力学は波動関数を抽象ヒルベルト空間に置き、一般相対性理論は重力を共有連続体の曲率に置く。いずれの構成も、単一の質量粒子が運ぶ代数から、物質波が\emph{何であるか}を述べない。

二重時空理論は候補を与える~\cite{dst-main}。すべての質量粒子は、16実次元双四元数代数 \(\mathbb{H}\otimes\mathbb{H}\cong\mathrm{Cl}(3,1)\) の内部に二つのコンパクト化ミンコフスキー枠を運ぶ。重力と慣性は通常ローター \(R_{\mathrm{usual}}\) と双対ローター \(R_{\mathrm{dual}}\) のねじれ不整合である。双対ローターは \(\mathbb{C}^2\) 上で
\[
R(\beta)=\exp\Bigl(\sum_{a=1}^3\tfrac{\beta_a}{2}(-i\sigma_a)\Bigr)
\]
として表され、ユニタリかつ行列式 \(1\)、したがって \(\mathrm{SU}(2)\) の元である。双対セクターのコンパクト性は、その群の \(4\pi\) 周期である。

問いは、\(R(\beta)\) のどの不変量がド・ブロイ位相であるか、である。指標、スピノル行列要素、不整合角から作ったスカラー \(\mathrm{U}(1)\) 因子は、三つの異なる量である。それらを混同すると、実のトレースから複素波動関数が現れ、あるいは振動しないラグランジアンからコンプトン振動数が現れる。三つを別々に名指したとき、幾何は明瞭になる。

\section{双対ローター}

質量粒子の完全ローターは
\[
R_{\mathrm{total}}
  =\exp\Bigl[\sum_{a=1}^3\bigl(\tfrac{\theta_a}{2}i\Gamma_a+\tfrac{\phi_a}{2}\Gamma_a\bigr)\Bigr],
\qquad \Gamma_1=I,\;\Gamma_2=J,\;\Gamma_3=K
\]
と書かれる。ここで \(i\Gamma_a\) はブーストを生成し（平方は \(+1\)）、\(\Gamma_a\) は双対セクターの回転を生成する（平方は \(-1\)）。同軸の通常生成子と双対生成子は可換であり、一軸の指数は因数分解する。異軸の生成子は必ずしも可換でない。無制限の六パラメータ因数分解 \(R_{\mathrm{usual}}R_{\mathrm{dual}}\) は、それらの交換子によっては正当化されない。

相対ローター \(\Omega=R_{\mathrm{usual}}^{\dagger}R_{\mathrm{dual}}\) が単位元に等しいのは、二つのローターが一致するときに限る。パラメータ空間スカラー \(J\) の消滅は \(\Omega=1\) を強制しない。平衡した質量シードは \(J=0\) でありながら非自明な相対ローターをもつ。

\(\mathbb{C}^2\) 上で双対ローターはロドリゲス行列
\[
R(\beta)=\cos(\lVert\beta\rVert/2)\,I+\sin(\lVert\beta\rVert/2)\,\widehat{\Gamma}(\beta)
\]
であり、\(\beta\neq 0\) のとき \(\widehat{\Gamma}(\beta)^2=-1\) である。通常セクターのローターは双曲の双子である。エルミートで行列式 \(1\)、非零ブーストに対してはユニタリでない。

\section{三つの不変量}

双対ローターの指標は正規化トレース
\[
\frac12\operatorname{Tr} R(\beta)=\cos(\lVert\beta\rVert/2)
\]
である。それは実であり、実であることに幾何的な理由がある。\(\lvert\uparrow\rangle\) と \(\lvert\downarrow\rangle\) を、観測者の双対枠における \(\sigma_z\) の計算基底固有ベクトルとする。トレースはそれらの和
\[
\operatorname{Tr} R(\beta)
  =\langle\uparrow\rvert R(\beta)\lvert\uparrow\rangle
  +\langle\downarrow\rvert R(\beta)\lvert\downarrow\rangle
\]
であり、したがって半トレースはスピンアップ行列要素の実部である。観測者軸に沿っては、二つの振幅は反対の半角位相
\[
\langle\uparrow\rvert R(\theta e_z)\lvert\uparrow\rangle=e^{-i\theta/2},\qquad
\langle\downarrow\rvert R(\theta e_z)\lvert\downarrow\rangle=e^{+i\theta/2}
\]
であり、その和は \(2\cos(\theta/2)\) である。このトレースから作った対は、複素スカラー波動関数ではありえない。二つの \(\mathrm{U}(1)\) 位相はすでに加えられ、虚部は相殺している。残るのは双対ラピディティの \(\mathrm{SU}(2)\) 指標であり、単位双対ローターの三次元球面上の回転の半角余弦である。

より精密な不変量は単一のスピノル行列要素である。\(\theta=\lVert\beta\rVert\)、\(\hat n=\beta/\theta\) として
\[
\langle\uparrow\rvert R(\beta)\lvert\uparrow\rangle
  =\cos(\theta/2)-i n_z\sin(\theta/2)
\]
である。その絶対値の二乗は、観測者に対する双対軸の極角傾斜
\[
\bigl\lvert\langle\uparrow\rvert R(\beta)\lvert\uparrow\rangle\bigr\rvert^2
  =\cos^2(\theta/2)+n_z^2\sin^2(\theta/2)
  =1-(n_x^2+n_y^2)\sin^2(\theta/2)
\]
である。単位絶対値の純位相が現れるのは、その軸が観測者に揃っているとき、あるいは回転が自明なときに限る。二つの軸が同じ幾何を明瞭に示す。観測者軸のまわりの双対回転は、すでに書いた位相 \(e^{-i\theta/2}\) である。赤道面内の双対回転は実余弦
\[
\langle\uparrow\rvert R(\theta e_x)\lvert\uparrow\rangle=\cos(\theta/2)
\]
である。したがって行列要素の振幅と位相は、単一の角の独立した装飾ではない。それらは双対軸が観測者に対してどう座っているかを記録する。教科書的波動力学の極形式 \(\sqrt{\rho}\,e^{iS/\hbar}\) はこの \(\mathrm{SU}(2)\) の対象ではなく、二つの実単位ベクトルに対して \(0\) または \(\pi\) でしかない \(\arg(\hat n_{\mathcal{O}}\cdot\hat n_A)\) から回復することもできない。

第三の量は、通常ラピディティと双対ラピディティのあいだのスカラー不整合角 \(\delta\) であり、手で実験室の \(\mathrm{U}(1)\) 因子 \(e^{i\delta}\) へ昇格される。その同定はド・ブロイ位相 \((Et-\mathbf{p}\cdot\mathbf{x})/\hbar\) として読める。それは指標ではなく、自動的にスピノル行列要素でもない。\(\mathrm{SU}(2)\) の二重被覆は区別を行列の恒等式にする。観測者軸のまわりの \(2\pi\) 双対回転は \(-I\) であり、スピノルはその負へ、半トレースは \(\cos\pi=-1\) へ送られる。\(4\pi\) 回転は単位元である。通常のスカラー物質波は、すでに \(2\pi\) で一周期を完了しているであろう。

\section{観測者重なり}

\(\mathcal{O}\) を巨視的観測者、\(A\) を質量粒子とする。双対ヒルベルト空間に住む比較は、参照スピノル \(\chi\) に対し
\[
\langle R(\beta_{\mathcal{O}})\chi\mid R(\beta_A)\chi\rangle
  =\langle\chi\mid R(\beta_{\mathcal{O}})^{\dagger}R(\beta_A)\chi\rangle
\]
である。各双対ローターのユニタリ性により、左辺は \(\chi\) の二つのユニタリ輸送された複製の内積であり、右辺は相対双対ローター \(R(\beta_{\mathcal{O}})^{\dagger}R(\beta_A)\) の行列要素である。その相対ローターはそれ自身ユニタリかつ行列式 \(1\) であり、対は \(\mathrm{SU}(2)\) を離れない。絶対値は \(\lVert\chi\rVert^2\) で抑えられる。

この対はパラメータ空間のキリング形式ではない。キリング対は符号 \((3,3)\) をもつ。自己重なりは \(J\) の定数倍に等しく、負になりうる。ボルン確率ではない。双対時空論理の四つのラベルは、\(\mathbb{C}^2\) の四本の直交光線ではない。

双対ローターの逆は反対ラピディティの双対ローターである：\(R(\beta)^{\dagger}=R(-\beta)\)。したがって相対ローターはそれ自身双対ローターの積 \(R(-\beta_{\mathcal{O}})R(\beta_A)\) である。共通の光線上では生成子は可換であり、積は尺度の差による単一の双対回転へ潰れる：\(R(s\hat n)^{\dagger}R(t\hat n)=R((t-s)\hat n)\)。それは共有軸ですでに見える角の加法を、輸送ではなく比較として読んだものである。

その軸を外れるとラピディティは加えない。四元数表 \(\Gamma_a\Gamma_b=-\delta_{ab}+\varepsilon_{abc}\Gamma_c\) はロドリゲス積をスカラーとトレース零の残余へ展開する。スカラーは合成の指標
\[
\tfrac12\operatorname{Tr}\bigl(R(\beta)R(\gamma)\bigr)
  =\cos(\theta/2)\cos(\phi/2)-(\hat n\cdot\hat m)\sin(\theta/2)\sin(\phi/2)
\]
であり、\(\theta=\lVert\beta\rVert\)、\(\phi=\lVert\gamma\rVert\)、\(\hat n\)、\(\hat m\) は二つの双対軸である。したがって積の指標は指標の積ではない。軸の内積によってずれる。相対ローターでは逆がその符号を反転し、
\[
\tfrac12\operatorname{Tr}\bigl(R(\beta_{\mathcal{O}})^{\dagger}R(\beta_A)\bigr)
  =\cos(\theta_{\mathcal{O}}/2)\cos(\theta_A/2)+(\hat n_{\mathcal{O}}\cdot\hat n_A)\sin(\theta_{\mathcal{O}}/2)\sin(\theta_A/2)
\]
となる。整列した軸は差の半角余弦を回復する。直交する軸は余弦の積だけを残す。単一の計算基底スピノルの対はなお行列要素であり、この指標ではない。指標はすでに二つの反対位相を加えている。

積の残余はトレース零である。それは \(\Gamma(\hat n)\)、\(\Gamma(\hat m)\)、および \(\Gamma(\hat n\times\hat m)\) の実線形結合である。外積が非可換な残余である。合成の軸を傾け、半トレースには見えない。したがって異なる双対軸の重ね合わせは \(\mathrm{SU}(2)\) における合成であり、複素スカラーの加法ではない。鋭い実例は、直交する軸のまわりの一対の \(\pi\) 双対回転である。それらは第三軸のまわりの \(\pi\) 回転へ合成し、順序を反対にすれば符号も反対になり、群交換子は \(-I\) である。観測者軸のまわりの \(2\pi\) 双対回転と同じ中心的反転である。この幾何における非可換干渉は、スカラー縞への小さな補正ではない。\(\mathrm{SU}(2)\) の二重被覆そのものである。

\section{不整合の動態}

六つのラピディティに対する粒子内在作用は模型であり、ユニタリ性からの演繹ではない。近傍の三つの系を比べる~\cite{dst-dynamic-equation,dst-main}。

書かれた反対符号の密度
\[
L=\tfrac12\dot\phi^2-\tfrac12\dot\theta^2-\tfrac m2(\phi-\theta)^2
\]
のオイラー–ラグランジュ方程式は \(\ddot\phi=-m(\phi-\theta)\)、\(\ddot\theta=-m(\phi-\theta)\) である。不整合 \(\delta=\phi-\theta\) は自由である：\(\ddot\delta=0\)。アフィン遅れ \(\delta(t)=\delta_0+\nu t\) は許される。係数 \(\nu\) は初期条件である。質量 \(m\) は両角をともに加速し、ド・ブロイ振動数を固定しない。

しばしば並記される近傍の系
\[
\ddot\phi=m(\phi-\theta),\qquad -\ddot\theta=m(\phi-\theta)
\]
は不整合を暴走させ、\(\ddot\delta=2m\delta\) とする。振動子ではない。

振動子が現れるのは、両運動項の符号を正にとったときに限る。密度 \(L=\tfrac12\dot\phi^2+\tfrac12\dot\theta^2-\tfrac m2(\phi-\theta)^2\) は \(\ddot\delta+2m\delta=0\) を与え、単位 \(\hbar=c=1\) で振動数 \(\sqrt{2m}\) である。その振動数はコンプトン振動数 \(m\) ではなく、模型は書かれた作用でもない。

したがって自由不整合の線形蓄積は、書かれた作用が実際に供給する、永年位相への唯一の動力学的橋である。その速さを \(E/\hbar\) と同定することは読みのままである。粒子内在的 \(J_{\mathrm{osc}}\) の多数粒子にわたる集団平均も同様に読みである。ワイツェンベック密度 \(T\) との同一視ではなく、\(\lvert J_{\rm norm}\rvert\le 1\) をここへ移植するのでもない。

コンプトン波長は、双対セクターのコンパクト化半径 \(\sim\hbar/mc\) として読み続けられる。ド・ブロイ波長は、同定されたスカラー位相が \(2\pi\) 進む実験室距離である。いずれの長さも、トレース指標から導かれるのではない。

\section{一次形式としてのディラック方程式}

基礎代数は \(\mathrm{Cl}(3,1)\) である。したがって一階のスピノル方程式は追加の要請ではなく、積の内部にすでに存在する平方根である。四元運動量 \(p\) に対し、次数1の元 \(\gamma(p)=\sum_{\mu}p_{\mu}\gamma^{\mu}\) は、代数においても \(\mathbb{C}^4\) 上のワイル行列としても \(\gamma(p)^2=Q(p)\) を満たす。ここで \(Q(p)=-p_0^2+p_1^2+p_2^2+p_3^2\) である。質量殻 \(Q(p)=-m^2\) は \((\gamma(p)-im)(\gamma(p)+im)\) の消滅である。\(\gamma(p)\Psi=im\Psi\) を満たすスピノルはその殻の上にある。カイラル表示ではスラッシュはブロック反対角であり、質量が零なら二つのワイルブロックは分離し、質量が非零ならそれらは結合する。カイラリティ \(\gamma^5\) は \(m\) の符号を反転する。静止系の解は存在する。静止系の粒子スピノルは、双対スピノル \(u\) からワイル対 \((u,\,iu)\) として組み立てられる。\(u\) の双対回転はその静止系方程式を保つ。したがって双対セクターの内部 \(\mathrm{SU}(2)\) は \(p=(m,0,0,0)\) における質量殻の対称性である。二つのワイルブロックは \(\mathbb{C}^4\) で直交し、双対回転だけ異なる二つの静止系解の内積は、双対ローター振幅の二倍である。
\[
\langle\Psi(u)\mid\Psi\bigl(R(\beta)u\bigr)\rangle_{\mathbb{C}^4}
  =2\,\langle u\mid R(\beta)u\rangle_{\mathbb{C}^2}.
\]
したがって双対ローターのスピノル行列要素は、ディラック方程式の傍らに置かれた余分な対象ではない。内部双対回転で結ばれた静止系解のヒルベルト対である。その対はなお \(e^{iEt/\hbar}\) ではなく、時空の偏微分方程式を与えるのでもない。

これが運動量空間の自由ディラック方程式であり、\((i\gamma^{\mu}\partial_{\mu}-m)\Psi=0\) の平面波形に同値である。スラッシュを平方するとクライン–ゴルドン記号 \(Q(p)+m^2\) が回復する。一階構造はクリフォード積から継承されるのであり、ねじれ真空のまわりでの双対ローター偏微分方程式の線形化からではない。\(m\) を双対ローター剛性と同一視することは幾何的な読みであり、導かれた等式ではない。時間ガンマ \(\gamma^0\) はミンコフスキーベクトル \(e_0\) であり、スピン射影に用いる双曲バイベクトルではない。

\section{波動方程式と幾何的な読み}

ゆっくり変化する双対ローター背景は、よく知られた実験室波動方程式
\begin{align}
i\hbar\frac{\partial\psi}{\partial t}
  &=\Bigl[-\frac{\hbar^2}{2m}\nabla^2+V(\mathbf{x})\Bigr]\psi, \\
\square\psi+\frac{m^2c^2}{\hbar^2}\psi&=0
\end{align}
を示唆するものとして読める。\(\mathrm{SU}(2)\) 運動学、あるいはアフィン不整合から、これらのスカラー偏微分方程式への移行は発見的展開であり、導かれた定理ではない。非相対論的および相対論的量子力学は、双対ローター動態の低エネルギー言語として利用可能であり続ける。それらは指標からもスピノル行列要素からも強制されない。

同じ調子で、波動関数の収縮は、多数の双対ローターが剛な集団状態にロックされた巨視的観測者へ、粒子の双対ローターが強制整列することとして読める。見かけのランダム性は、粒子あたり六つの実数である初期双対角に対する無知である。双対セクターはコンパクトであるため、再整列は双対次元において瞬時である。これは測定要請の幾何的描像であり、その導出ではない。

同一の16次元ネットワークは、平均された不整合を通じて巨視的重力を、局所的双対ローター幾何を通じて微視的干渉を伝えるものとして読める。本稿の診断量 \(J\) は主論文のパラメータ空間キリング二次形式ではない。重力と量子非局所性は異なるスケールにおける同一の代数でありうる。それは読みのままである。

\section{ド・ブロイ–ボーム力学との関係}

ド・ブロイとボームのパイロット波解釈は、粒子を決定論的軌跡に沿って操舵する外部誘導場を仮定する~\cite{bohm1952}。二重時空枠組みには外部場は不要である。観測者重なりに現れる双対ローターは、粒子自身の双対ローターを重心の仮想位置で評価したものである。隠れた変数という言葉を用いるなら、それは六つの実ローター角である。「パイロット波」は双対セクターから見た粒子である。

ボーム力学はなお、時空上のスカラー波動関数と量子ポテンシャルを必要とする。双対ローター幾何は \(\mathrm{SU}(2)\) 運動学と相対ローター対を供給する。それ自体で量子ポテンシャルを再構成するのでもなく、ベル不等式の破れを証明するのでもない~\cite{bell1964}。双対セクターのコンパクト性は、非局所性を幾何的な選択肢として利用可能にするのであり、指標の定理としてではない。

\section{実験的兆候}

読みとして残る同定は、なお区別可能な期待を与える。

第一に、物質波干渉法は、近傍質量の集団的ねじれ剛性に対する干渉縞可視度の弱い依存と、大きな回転質量の角速度に線形な位相ずれ——双対ローターのフレームドラッグから生じる重力ファラデー様効果——を見るであろう。

第二に、コンプトン振動数近傍で作動する高周波重力波検出器が共鳴特徴を見るのは、調和的双対ローター模型が物理的に実現されている場合に限る。作業用振動子が正しい模型であるなら、それは \(\sqrt{2m}\) にあり \(m\) にはない。いずれの振動数における細線の不在も動態を制約するのであり、双対ローター運動学を棄却するのではない。

第三に、収縮が双対ローター同期であるなら、ローター緩和時間内の反復弱測定はボルン則からの微小な系統的ずれを示すであろう。この種の実験は、超伝導量子ビットまたは捕捉イオンの到達範囲内にある。

これらの兆候は、物質波の二重時空起源を、標準量子力学および外部パイロット場を保持する他の隠れた変数理論から区別する。それらは代数的である指標の恒等式やスピノル行列要素を検定するのではない。

\section{結論}

物質波は双対ローターのトレースではない。そのトレースが実であるのは、二つの反対の半角位相を足すからである。双対セクターが実際に供給するのは \(\mathbb{C}^2\) 上の \(\mathrm{SU}(2)\) 運動学である。実の指標、絶対値が双対軸の観測者に対する極角傾斜である複素スピノル行列要素、そしてその群を離れずに二つの双対ローターを比べる相対ローター対。軸外では対の指標が軸の整列を記録し、非可換な残余はトレース零である。同じ行列要素は、内部双対回転だけ異なる静止系ディラック解の重なりである。スカラーのド・ブロイ位相はさらなる同定であり、書かれた粒子作用における線形遅れとして利用可能であって、静止質量から導かれた調和運動としてではなく、\(\mathrm{Cl}(3,1)\) にすでに含まれる一次のディラック因子としてでもない。

この幾何における量子力学は、重力の傍らに置かれた余分なヒルベルト空間ではない。同一の16次元代数の双対ローター運動学である。運動学は鋭い。実験室波動方程式への動力学的橋は、そうではない。

\begin{thebibliography}{9}
\raggedright

\bibitem{dst-main}
Anonymous,
「双対時空の間のねじれとしての重力：一般相対性理論の双四元数的再定式化」 (2025).
\url{https://github.com/hypernumbernet/dual-spacetime-doc/blob/main/double-spacetime-theory.pdf}.

\bibitem{dst-dynamic-equation}
Anonymous,
「二重時空理論における双対ローター動態：書かれた作用、自由不整合、および近傍の振動子」 (2025).
\url{https://github.com/hypernumbernet/dual-spacetime-doc/blob/main/dst-dynamic-equation.pdf}.

\bibitem{debroglie1924}
L.~de Broglie,
``Recherches sur la th\'eorie des quanta,''
\textit{Ann.\ Phys.} (Paris) \textbf{3}, 22--128 (1925).

\bibitem{bohm1952}
D.~Bohm,
``A suggested interpretation of the quantum theory in terms of ``hidden'' variables.\ I,''
\textit{Phys.\ Rev.} \textbf{85}, 166--179 (1952).

\bibitem{bell1964}
J.~S.~Bell,
``On the Einstein Podolsky Rosen paradox,''
\textit{Physics Physique Fizika} \textbf{1}, 195--200 (1964).

\end{thebibliography}

\end{document}
